\chapter{Risk Assessment}
\label{chap:risk_assessment}

\section{\ac{RD} Department}

The \ac{RD} Department at GreenGrid Energy is dedicated to innovation and the advancement of smart grid technologies. Given its critical role in safeguarding intellectual property and proprietary technologies, it is essential to systematically identify, assess, and mitigate risks related to cybersecurity threats, data integrity, and operational disruptions.

This risk assessment aligns with the ISO/IEC 27001 standard and forms part of the organization's \ac{ISMS}.

\subsection{Asset Inventory}

Key assets within the scope of this assessment include:
\begin{itemize}
  \item \textbf{Prototype Simulation Clusters:} Essential for testing and validating new technologies.
  \item \textbf{Algorithm and Software Development Platforms:} Contain proprietary algorithms and code.
  \item \textbf{Research Databases and Cloud Repositories:} Storage and processing of experimental data and intellectual property.
\end{itemize}

These assets have been classified based on their \ac{CIA} requirements:

\begin{tabular}{|l|c|c|c|}
\hline
\textbf{Asset} & \textbf{Confidentiality} & \textbf{Integrity} & \textbf{Availability} \\
\hline
Prototype Simulation Clusters & High & High & Medium \\
Algorithms and Software Platforms & High & High & High \\
Research Databases and Repositories & High & High & High \\
\hline
\end{tabular}

\subsection{Risk Identification}

Risks specific to the \ac{RD} environment were identified as follows:
\begin{itemize}
  \item \textbf{Data Breach or Intellectual Property Theft:} Unauthorized external or internal access.
  \item \textbf{Cyberattacks on Testing Environments:} Potential compromise through ransomware or malware.
  \item \textbf{Access Control Vulnerabilities:} Inadequate management of user privileges.
  \item \textbf{Data Loss or Corruption:} Insufficient backup and disaster recovery procedures.
\end{itemize}

\subsection{Risk Analysis}

Using GreenGrid’s Risk Level Matrix, the risks were analyzed based on their likelihood and potential impact:

\begin{tabular}{|p{9cm}|c|c|c|}
\hline
\textbf{Identified Risk} & \textbf{Likelihood} & \textbf{Impact} & \textbf{Risk Level} \\
\hline
Unauthorized access to proprietary IP & Likely & Major & High \\
Ransomware targeting \ac{RD} systems & Possible & Critical & High \\
Insider threat leading to data compromise & Possible & Major & High \\
Data corruption in testing environments & Unlikely & Moderate & Medium \\
Lack of \ac{MFA} & Likely & Moderate & Medium \\
\hline
\end{tabular}

\subsection{Risk Assessment and Treatment}

Risks categorized as High or Medium require defined treatment strategies:

\paragraph{High Risks:}
\begin{itemize}
  \item \textbf{Unauthorized access to proprietary IP:} Implement strict \ac{RBAC}, data encryption, and audit logging.
  \item \textbf{Ransomware targeting \ac{RD} systems:} Enhance endpoint security, implement antivirus and antimalware solutions, and perform continuous vulnerability assessments.
  \item \textbf{Insider threats:} Establish comprehensive access reviews, enhanced logging, and regular employee security training.
\end{itemize}

\paragraph{Medium Risks:}
\begin{itemize}
  \item \textbf{Data corruption in testing environments:} Strengthen backup procedures and validate backup integrity regularly.
  \item \textbf{Lack of \ac{MFA}:} Enforce \ac{MFA} across all \ac{RD} platforms to prevent unauthorized access.
\end{itemize}

\subsection{Implementation and Monitoring}

The \ac{RD} department will implement the following:
\begin{itemize}
  \item \textbf{Encryption} for sensitive data, both at rest and in transit.
  \item \textbf{Quarterly audits} of user access rights and system vulnerabilities.
  \item \textbf{\ac{IRP}} tailored to the \ac{RD} environment.
  \item \textbf{Regular employee training and awareness programs} focusing on cybersecurity best practices.
\end{itemize}
\chapter{Risk Assessment Report}
\label{chap:risk_assessment_report}

This chapter presents the results of the risk assessment for GreenGrid Energy. 
The assessment is in accordance with ISO/IEC 27001 Clause 8.2, using the risk management methodology defined in Chapter~\ref{chap:risk_management} and based on the critical assets documented in Appendix~\ref{appendix:asset_inventory}.

The assessment process followed the GreenGrid Energy risk matrix presented in Figure~\ref{figure:risk_matrix} and prioritised threats and vulnerabilities in terms of likelihood and impact.

\section{Risk Assessment Results}
The following subsections present the threats, vulnerabilities and assessed risks associated with common assets and each specific departmental asset, as identified in the asset inventory and assessed using the defined methodology.

\subsection{Shared Assets}

\begin{longtable}{|p{2cm}|p{2.5cm}|p{3cm}|p{2cm}|p{2cm}|p{2cm}|}
\hline
\textbf{Asset ID} & \textbf{Threat} & \textbf{Vulnerability} & \textbf{Likelihood} & \textbf{Impact} & \textbf{Risk Level} \\
\hline
\endhead
ERP-301 & Data Breach & Poor access controls & Possible & Severe & \textbf{Medium-High} \\
\hline
IT-201 & Phishing Attacks & Insufficient staff training & Very Likely & Significant & \textbf{High} \\
\hline
IT-202 & Unauthorised Access & Firewall misconfiguration & Possible & Significant & \textbf{Medium-High} \\
\hline
CS-401 & Data Leakage & Improper permission management & Possible & Significant & \textbf{Medium-High} \\
\hline
HR-501 & Insider Threat & Excessive user privileges & Possible & Significant & \textbf{Medium-High} \\
\hline
HR-502 & Staff Unavailability & High dependency on key staff & Possible & Significant & \textbf{Med-High} \\
\hline
\end{longtable}

\subsection{Supply Chain Management Assets}

\begin{longtable}{|p{2cm}|p{2.5cm}|p{3cm}|p{2cm}|p{2cm}|p{2cm}|}
\hline
\textbf{Asset ID} & \textbf{Threat} & \textbf{Vulnerability} & \textbf{Likelihood} & \textbf{Impact} & \textbf{Risk Level} \\
\hline
\endhead
ERP-301 & Supply Chain Compromise & Insecure third-party integrations & Possible & Severe & \textbf{Medium-High} \\
\hline
ERP-301 & Data Integrity Loss & Configuration errors, incorrect updates & Possible & Significant & \textbf{Medium-High} \\
\hline
ERP-301 & System Downtime & Single point of failure & Unlikely & Severe & \textbf{Medium-High} \\
\hline
\end{longtable}

\section{Recommended Mitigation Actions and Control Mapping}
According to the criteria defined in chapter~\ref{chap:risk_management}, all risks identified in this assessment are classified as \textbf{High} or \textbf{Medium-High} and therefore need to be addressed by treatment strategies.
The following table maps each risk to relevant ISO/IEC 27001:2022 Annex A\cite{ismsonline2025annexa} controls and provides a justification for inclusion.

\begin{longtable}{|p{2.2cm}|p{4.5cm}|p{4.5cm}|p{4cm}|}
\hline
\textbf{Asset ID} & \textbf{Threat} & \textbf{Control(s) (Annex A)} & \textbf{Justification} \\
\hline
\endhead
ERP-301 & Data Breach (Access Control Failure) & A.9.1.1, A.9.4.1 & Access control policies must be enforced. Exclusion not justified due to system criticality. \\
\hline
IT-201 & Phishing / Social Engineering & A.6.3.1, A.7.2.2 & Mandatory training and awareness campaigns are needed. Control inclusion is essential. \\
\hline
IT-202 & Network Misconfiguration & A.13.1.1, A.12.4.1 & Logging and network security controls reduce the risk of unauthorised access. Must be included. \\
\hline
CS-401 & Data Leakage & A.9.2.3, A.9.2.5 & Access rights and permissions must be reviewed periodically. Exclusion unjustified. \\
\hline
HR-501 & Insider Threat & A.9.2.5, A.12.4.3 & Oversight of privileged accounts and logging of admin actions are critical. Controls mandatory. \\
\hline
HR-502 & Staff Unavailability & A.17.1.1, A.5.30 & Continuity plans and redundancy for key roles must be established. Justified and required. \\
\hline
ERP-301 & Supply Chain Compromise & A.15.1.1, A.15.2.1 & Supplier vetting and security requirements must be formalised. Cannot be excluded. \\
\hline
ERP-301 & Data Integrity Loss & A.12.1.2, A.14.2.2 & Change and update management controls are required. \\
\hline
ERP-301 & System Downtime & A.17.1.2, A.17.2.1 & Backup and failover mechanisms must be in place. Control is mandatory for operational continuity. \\
\hline
\end{longtable}

\section{Implementation Considerations}

All risks identified in this assessment fall into the \textbf{High} or \textbf{Medium-High} categories and according with risk appetite defined on section \ref{subsubsection:risk_apetite} require immediate attention.  Each control has been explicitly mapped to ISO/IEC 27001 Annex A and none have been excluded - each risk affects a core system with significant operational or compliance impact.

In order to translate these recommendations into effective actions, GreenGrid Energy should establish a control implementation plan that:

\begin{itemize}
  \item \textbf{Assigns clear ownership:} Identify the person or team responsible for implementing, maintaining and periodically reviewing each control.
  \item \textbf{Defines timing and resources:} Establish realistic timelines, budget allocations, staffing requirements and any necessary tools or training.
  \item \textbf{Defines monitoring and metrics:} Identify key performance indicators and the mechanisms by which effectiveness will be tracked.
  \item \textbf{Aligns with the ISMS:} Ensure that all activities integrate seamlessly with the organisation's wider \ac{ISMS} processes and the continuous improvement cycle outlined in chapter~\ref{chap:risk_management}.
\end{itemize}

Once controls have been implemented, each risk addressed shall be reassessed as part of the ongoing monitoring and review of the \ac{ISMS} to confirm that the residual risk remains within acceptable limits and that lessons learned are incorporated into future assessments.
